% Physical pages 3 and 4 of the concise study: the movement of the formulary,
% stated as an argument. The developed comparison of the three interpretations
% begins on the fifth page and deepens this account without repeating it.

\section{The Propers: Themes and Movement}
\zlabel{triptych:concise:themes:start}

This Mass asks who comes to God's feast, and what keeps a person from it. A
king makes a wedding for his son and sends to call the invited, ``and they
would not come''; the hall is filled instead from the roads with whomever the
servants find, ``both bad and good'' (Mt 22:3, 10). Isaiah has already spread
the table: on his mountain the Lord makes ``unto all people'' a feast of fat
things and of wine, and casts death down for ever (Is 25:6, 8). Around these
two readings the other texts set a confession that no one stands on his own
account, a prayer for grace at both ends of every good work, an apostle held
neither by plenty nor by want, and a table at which the communicants ask to
share the divine nature.

\begin{itemize}[leftmargin=1.2em,itemsep=0pt,topsep=.3em]
\item \textbf{Grace before the call:} the Entrance Antiphon and the Collect.
\item \textbf{The feast promised and the table spread:} the first reading and
the psalm.
\item \textbf{Freedom and hope:} the second reading and the Alleluia verse.
\item \textbf{The king's wedding:} the Gospel, in its longer or shorter form.
\item \textbf{Offering, table and vision:} the Prayer over the Offerings, the
Communion antiphons and the Prayer after Communion.
\end{itemize}

The antiphons and orations are the Missal's for the Twenty-eighth Week in
Ordinary Time and serve all three years; the readings, psalm and acclamation
are Year~A's, Lectionary no.~142. The books state their own relations: the
first reading is chosen for its relation to the Gospel, the psalm responds to
it, and the acclamation greets the Gospel. Philippians, read on its fourth and
last Sunday, and Matthew each follow a course of their own.

The \work{Ordo lectionum Missae} sets a title over each reading:
\latin{Faciet Dominus convivium, et absterget lacrimam ab omni facie}, the Lord
will make a feast and wipe away the tear from every face, from Isaiah 25:6 and
8; \latin{Quoscumque inveneritis, vocate ad nuptias}, whomever you find, call
to the wedding, from Matthew 22:9, a verse both forms of the Gospel contain;
and \latin{Omnia possum in eo, qui me confortat}, I can do all things in him
who strengthens me, from Philippians 4:13. The two titles that belong together
set a feast for every face beside a call to whoever is found.

\subsection{A question no one can answer alone}

The Entrance Antiphon is made from Psalm 129, a song of ascents. It asks who
could stand if the Lord marked iniquities, and answers that with him there is
propitiation; it ends by naming him the God of Israel. The Collect, \latin{Tua
nos, quaesumus, Domine, gratia}, asks that his grace \latin{semper et
praeveniat et sequatur}, always go before and follow, and keep the assembly
constantly intent on good works. Before any invitation is heard, the Mass has
asked for grace at the beginning and at the end of what it does.

\subsection{A feast on the mountain, a table in the house}

Isaiah's feast follows a song of thanks to the God who has been ``a strength
to the poor'' (25:4). On ``this mountain'' the Lord of hosts makes the feast
for all people, destroys the bond that held them, casts death down headlong,
wipes away tears from every face and takes away his people's reproach; and in
that day they say, ``Lo, this is our God, we have waited for him'' (25:9).

The psalm answers in the first person. The Lord rules the singer, who shall
want nothing; the shepherd's images give way to a host's, a table prepared in
the sight of enemies, the head anointed, the cup filled, and mercy following
the guest into the house of the Lord. What Isaiah promises to all peoples, the
psalm receives as one guest's table.

\subsection{Full and hungry, and the hope of the calling}

The second reading comes from the close of Philippians. Paul has learned to be brought low and to abound,
``to be full and to be hungry'', and he can ``do all things in him who
strengtheneth me'' (4:12--13). He commends the Philippians for sharing his
tribulation; the verses on their gift are left out, and the reading resumes
with God's supplying of all their want and a doxology. That supplying is a
wish in the Vulgate and a promise in the \work{Nova Vulgata}; both name what is
filled \latin{desiderium}, desire.

The Alleluia verse turns Paul's prayer for the Ephesians into the assembly's
prayer for itself: that the Father enlighten the eyes of our heart, to know
the hope of our calling. The word ``calling'' stands in the verse that greets
a Gospel about those who were called.

\subsection{The king's wedding}

The parable is the third Jesus tells in the temple to the chief priests and
the ancients of the people (21:23), after the vineyard parable's saying that
the kingdom shall be taken from them (21:43), at which the chief priests and
Pharisees ``knew that he spoke of them'' (21:45). The king's invitation is
refused twice, the second time
though ``all things are ready''. Some neglect it and go off, ``one to his farm
and another to his merchandise''; the rest seize the servants and kill them,
and the king's armies destroy the murderers and burn their city. The servants
are then sent to the places where the roads go out, \latin{ad exitus viarum},
and they gather bad and good until the hall is full.

In the longer form the king comes in to see the guests, finds one without a
wedding garment and asks him, ``Friend, how camest thou in hither not having
on a wedding garment?'' The man
is silent, and is bound and cast out, ``For many are called, but few are
chosen'' (22:11--14). The shorter form ends with the hall full. The choice
between the forms is pastoral, and nothing in the books makes either the
proper one; what rests on the garment and the closing saying holds only where
the longer form is read.

\subsection{Offering, table and vision}

The Prayer over the Offerings asks the Lord to receive the faithful's prayers
with the sacrificial gifts, that through these acts of devout worship they may
pass over to heavenly glory. At Communion the Missal offers either the psalm
verse on the rich who have known want and hunger while those who seek the
Lord lack no good thing, or John's promise that when the Lord appears his own
will be like him, because they will see him as he is. The Prayer after
Communion asks that those fed with the Body and Blood of Christ be made
partakers of the divine nature, the words of 2~Peter 1:4. The feast promised
on the mountain and refused in the parable ends as a table at which the
communicants ask to pass into glory, to see, and to share.

\subsection{Three questions, three readings}

The Fathers and Doctors who expounded these passages comment on single texts,
none on this Mass as a whole; each reading of the whole below is the editor's
synthesis, and every witness is credited only with what he says of his own
passage. Three questions divide them. The first asks
what the feast is and for whom it is made: Gregory the Great reads the
wedding as the Church joined to the Son, Jerome and Aquinas read Isaiah's
banquet of Christ, and Aquinas alone sets that banquet beside the king's
dinner. The second, which only the longer form raises, asks why a guest the
king's own servants gathered is cast out: Augustine and Gregory name the
missing garment charity, and other witnesses name it otherwise. The third asks
why the invited would not come: Gregory, Hilary, Chrysostom and Aquinas weigh
the farm and the merchandise. The three answers are compatible; the
witnesses part over what the garment is, over where the wedding stands, and
over the burnt city.
\zlabel{triptych:concise:themes:end}
